\begin{textsample}{NEG Dim 1 – 2015\_c5 – Score 7.00 – 2015\_c5/enade\_2015\_tec\_com\_ext\_a54eb9745d.txt}  \label{ex:f1_neg_348}
CAPÍTULO 5 DISTRIBUIÇÃO DOS CONCEITOS 5.1 PANORAMA NACIONAL DA DISTRIBUIÇÃO DOS CONCEITOS O Conceito Enade é um indicador de qualidade que avalia o desempenho dos estudantes a partir dos resultados do Enade.
Os valores possíveis vão de 1 ( pior situação ) a 5 ( melhor situação ).
As Unidades de Observação com apenas um ou sem nenhum concluinte participante não obtêm o Conceito Enade, ficando Sem Conceito ( SC ).
A Tabela 5.1 apresenta a quantidade e distribuição de cursos de Tecnologia em Comércio Exterior participantes do Enade/2015, por faixa de conceito e Grande Região.
Mais uma vez ressaltando que a diferença entre os cursos tabulados neste capítulo e no capítulo 2 corresponde aos cursos sem conceito, em princípio, aqueles sem alunos concluintes que participassem da prova ou que tivessem tão somente um aluno concluinte¹⁸.
Observando -se os dados da Tabela 5.1, nota -se que, dos 80 cursos participantes, 26 ( 32,5 \% ) classificaram -se com conceito 3, o valor modal.
Este foi também o conceito modal em mais três regiões : Sudeste ( 33,3 \% ), Sul ( 31,3 \% ) e Centro-Oeste ( 50,0 \% ).
A região Sul também apresentou valor modal ( 31,3 \% ) no conceito 4, caracterizando -se, portanto, nesta região, uma distribuição da moda em dois valores contíguos.
Já a região Centro-Oeste apresentou como valores modais os conceitos 3 e 1 ( 50 \% \textbf{cada} ).
Na região Nordeste, a moda foi o conceito 4 ( 66,7 \% ).
O conceito 4 foi o segundo mais frequente em nível nacional ( 26,3 \%, correspondendo a 21 cursos ) e o conceito 5, o terceiro ( 11,3 \%, correspondendo a nove cursos ).
Houve, ainda, oito cursos ( 10,0 \% ) que receberam conceito 2 e três cursos ( 3,8 \% ) que receberam conceito 1.
Dos 80 cursos de Tecnologia em Comércio Exterior, 13 cursos ( 16,3 \% ) ficaram sem conceito ( SC ). ¹⁸ Estes cursos com somente um concluinte não foram considerados no capítulo 2.
A região Norte participou com apenas dois cursos ou 2,5 \% do total nacional.
Nesta região, o conceito 5 foi atribuído a um dos dois cursos participantes, o que equivale a 50,0 \% do total regional.
Um curso ( 50,0 \% ) ficou sem conceito ( SC ), enquanto nenhum curso obteve os demais conceitos nessa região.
A região Nordeste participou com três cursos ou 3,8 \% do total nacional.
Nesta região, dois cursos, 66,7 \% em termos regionais, obtiveram conceito 4, o conceito modal para a região.
O conceito 3 foi atribuído a um único curso ( 33,3 \% ).
Nenhum curso ficou sem conceito ( SC ) ou recebeu os conceitos restantes, nessa região.
Dos 57 cursos participantes da região Sudeste ( 71,3 \% do total nacional ), o conceito 3 também foi o conceito modal, tendo sido obtido por 19 cursos ( 33,3 \% ).
O conceito 1 foi atribuído a dois cursos ( 3,5 \% ) e os conceitos 2 e 4, a seis e 14 cursos, correspondendo a 10,5 \% e 24,6 \%, respectivamente.
O conceito 5 foi atribuído a cinco cursos ( 8,8 \% ).
E 11 cursos ( 19,3 \% ) ficaram sem conceito ( SC ).
Os 16 cursos da região Sul ficaram distribuídos em quase todas as faixas de conceitos.
Dez cursos obtiveram conceitos 3 ou 4 ( cinco cursos \textbf{cada} ).
O conceito 2 foi atribuído a dois cursos ( 12,5 \% ), e o conceito 5, a três cursos ( 18,8 \% ).
Apenas um curso na região Sul ficou sem conceito ( SC ), e nenhum recebeu conceito 1.
Já dos dois cursos participantes na região Centro-Oeste, um recebeu conceito 1, enquanto o outro obteve conceito 3.
Nesta região, nenhum curso ficou sem conceito ( SC ) ou recebeu os demais conceitos.
O Gráfico 5.1 apresenta a distribuição do Conceito Enade, segundo Grande Região : poligonais mais à esquerda representam Grandes Regiões com pior distribuição, e poligonais mais à direita, Grandes Regiões com melhores conceitos.
A região Centro-Oeste ( linha preta ) apresenta os cursos em IES com conceitos mais baixos, ao passo que a região Norte ( linha azul ) apresenta a distribuição com valores maiores.
Cursos em IES das regiões Nordeste ( linha verde ), Sudeste ( linha vermelha ) e Sul ( linha roxa ) apresentam uma situação intermediária.
Deve -se levar em conta na análise que na região Norte, o único curso avaliado recebeu o conceito máximo. 5.2 CONCEITOS POR CATEGORIA ADMINISTRATIVA E POR GRANDE REGIÃO A Tabela 5.2 apresenta a distribuição dos cursos participantes do Enade/2015 de Tecnologia em Comércio Exterior, por Categoria Administrativa, de acordo com os conceitos por eles alcançados, segundo as Grandes Regiões brasileiras.
Dos 80 cursos participantes, oito ( 10,0 \% ) eram ministrados em Instituições Públicas e 72 ( 90,0 \% ), em Instituições Privadas.
De acordo com as informações da Tabela 5.2, em termos nacionais, entre Instituições Públicas e Privadas observa -se que, dos nove cursos que receberam conceito 5, um era proveniente de IES Pública, e oito, de IES Privadas.
Dos oito cursos participantes de IES Públicas, o conceito 4 foi o valor modal, atribuído a seis cursos ( 75,0 \% ).
Entre os demais cursos de Instituições Públicas participantes, um curso ( 12,5 \% ) foi avaliado com conceito 3, enquanto outro curso obteve conceito 5.
Nesta categoria, nenhum curso ficou sem conceito ( SC ) ou recebeu os conceitos 1 ou 2.
Na rede privada, o conceito modal foi 3, atribuído a 25 cursos dos 72 da categoria.
Entre os demais cursos participantes, três receberam conceito 1, e oito, conceito 2.
O conceito 4 foi atribuído a 15 cursos e o conceito 5, a oito cursos.
Nesta Categoria Administrativa, outros 13 cursos ficaram sem conceito ( SC ).
Na análise por região, observa -se que, na região Norte, as Instituições Privadas participaram com dois cursos ( o total regional ), dos quais um ficou sem conceito ( SC ), e outro obteve conceitos 5.
Nenhum curso recebeu os demais conceitos nessa combinação de categoria e região.
As Instituições Públicas não tiveram participação com nenhum curso na região Norte.
Na região Nordeste, a rede privada concentrou dois dos três cursos participantes, o equivalente a 66,7 \% do total da região.
Dentre os cursos oferecidos por IES Privadas no Nordeste, um curso ficou alocado ao conceito 3, outro curso obteve conceito 4.
Nenhum curso ficou sem conceito ( SC ) ou recebeu os demais conceitos restantes dessa categoria e região.
As Instituições Públicas dessa região participaram somente com um curso, alocado ao conceito 4.
Nenhum curso ficou sem conceito ( SC ) ou recebeu os conceitos restantes nesta combinação de categoria e região.
Na região Sudeste, a proporção de cursos da rede privada foi de 91,2 \%, correspondendo a 52 dos 57 cursos participantes.
Nesta categoria, na região Sudeste, o conceito modal foi 3 ( 19 cursos ).
Nesta combinação de Categoria Administrativa e Grande Região, 11 cursos ( 21,2 \% ) não receberam conceito ( SC ).
Os demais foram avaliados com conceito 1 ( dois cursos ), conceito 2 ( seis cursos ), conceito 4 ( dez cursos ) e conceito 5 ( quatro cursos ).
Entre os cinco cursos oferecidos por Instituições Públicas na região Sudeste, o conceito modal foi 4 ( quatro cursos ), e um curso recebeu o conceito 5.
Na região Sudeste, nesta categoria, nenhum curso ficou sem conceito ( SC ) ou recebeu os conceitos restantes.
As Instituições Privadas concentraram 15 dos 16 cursos participantes da região Sul, 93,8 \% do total regional.
Desses, cinco ficaram com conceito 3, o conceito modal.
Nesta combinação de Categoria Administrativa e Grande Região, quatro cursos ficaram com conceito 4, dois cursos obtiveram conceito 2, enquanto três cursos receberam conceito 5.
Nenhum curso foi associado ao conceito 1, enquanto um curso ficou sem conceito ( SC ).
As Instituições Públicas na região Sul participaram com somente um curso, e este foi alocado ao conceito 4.
Nenhum curso ficou sem conceito ( SC ) ou recebeu os demais conceitos restantes nessa categoria e região.
Na região Centro-Oeste, obteve -se dois cursos participantes, dos quais um deles era de Instituição Privada, e alocado ao conceito 1.
Ainda sobre esta categoria nota -se que nenhum curso ficou sem conceito ( SC ) ou recebeu os demais conceitos.
Também um único curso foi oferecido por Instituições Públicas, e este recebeu conceito 3.
Nesta região, nenhum curso de IES pública ficou sem conceito ( SC ) ou recebeu os conceitos restantes dessa categoria.
O Gráfico 5.2 apresenta a distribuição do Conceito Enade segundo a Categoria Administrativa da IES.
Os cursos em IES Públicas ( linha azul ) apresentam uma poligonal

% matched lemmas: cada
\end{textsample}
